\chapter{物质的三态远远不止三态}

\begin{kaichang}
夏天你网购了一盒冰淇淋，快递箱里除了冰淇淋，还有几个鼓鼓的袋子，里面装着几块冒着“白烟”的“冰”。你隔着袋子碰了碰——嘶，比冰箱里的冰块冷多了。说明书写着：干冰，请勿用手直接触碰。

几个小时后，冰淇淋吃完了，你想起那几块“冰”，打开袋子一看：空了。袋子里没有水，桌面是干的，垃圾桶里也没有任何残渣。那几块“冰”连一滴水都没留下，就这么消失了。

先别往下看，猜三个问题：它去了哪里？那团“白烟”是什么？为什么普通的冰会化成一摊水，它却干干净净地退场？这一章结束时，你要能理直气壮地回答这三问。
\end{kaichang}

\section{看得见的三态}

冰、水、水蒸气——这是你从小就知道的“物质三态”：固体、液体、气体。先不管为什么，只看现象，你能总结出什么？

固体有自己的形状和体积。一块冰放在碗里是方块，放在盘子里还是方块。液体有体积，但没有自己的形状：一杯水倒进碗里是碗的形状，倒进瓶子里是瓶子的形状。气体两样都没有：一瓶可乐打开后，跑出来的气体会占满整个房间，你既不能给它塑形，也没法把它“堆”在某个角落。

再看几个你多半见过的场面。烧开水时，壶嘴冒出一团白色的“气”，但仔细看，白气并不是贴着壶嘴出来的，而是离壶嘴一小段距离才出现。衣柜里的樟脑丸，过几个月变小了一圈，可柜子底下从来没有一摊“樟脑水”。冬天湿衣服晾在室外，先冻得硬邦邦，但过几天居然也干了。还有严寒早晨窗玻璃上的冰花——屋里并没有人往玻璃上泼水。从冰箱里拿出的饮料罐，外壁很快挂满水珠，罐子明明是密封的，水从哪里来？加油站旁的大罐上写着“液化石油气”——气体怎么能“液化”着存？这些场面里，物质都在改变“状态”，却似乎各有各的改法：有的经过液体，有的绕过液体，还有的干脆从空气里凭空现身。要把它们统一起来，得把镜头推进到肉眼看不见的地方。

\section{镜头推进：粒子在做什么}

第 1 章说过，物质由极小极小的粒子组成——水由水分子组成，冰也是，水蒸气也是。三种状态里，粒子本身没有任何不同，变的是它们的\textbf{排列方式}和\textbf{运动方式}。

\begin{itemize}[itemsep=2pt]
\item \textbf{固体}：粒子紧紧挨在一起，各自被邻居“卡”在一个固定位置上。它们并没有睡着，而是在原地不停振动，像一屋子人各自站在自己的位置上抖腿。位置固定，所以固体有固定的形状。
\item \textbf{液体}：粒子仍然紧挨着，但不再被卡死，可以互相挤着滑来滑去，像高峰期地铁车厢里的乘客——人挨着人，却能慢慢挪到门口。能滑动，所以液体能流动、没有固定形状；仍然紧挨，所以体积基本不变。
\item \textbf{气体}：粒子之间的距离一下子拉得很远，各自高速乱飞，撞来撞去，像被撒进大操场的一小撮弹珠。间距大、随便跑，所以气体既没有固定形状，也没有固定体积。
\end{itemize}

那么“温度”在这个图像里是什么？温度就是\textbf{粒子运动剧烈程度的标尺}。同样一杯水，温度高，分子平均就跑得快、撞得狠；温度低，就动得慢。你手摸到的“烫”，其实是大量高速粒子撞击你的皮肤、把能量交给你的感觉。

“压强”呢？气体粒子不停撞击容器壁，亿万次微小的撞击叠加起来，就是压强。给自行车胎打气，打进去的气越多，粒子撞壁越频繁，轮胎就越硬。把压扁的乒乓球泡进热水里，球内气体受热，粒子撞得更猛，就把凹陷顶了回去。温度、压强这些“宏观数字”，全都是粒子运动的集体表现。

\section{六个名字，一条规律}

有了粒子图像，那些“各有各的改法”就可以统一命名了。三种状态两两之间共有六个变化，你其实都见过：

\begin{itemize}[itemsep=2pt]
\item 固 $\rightarrow$ 液：\textbf{熔化}（冰化成水）；液 $\rightarrow$ 固：\textbf{凝固}（水冻成冰）。
\item 液 $\rightarrow$ 气：\textbf{汽化}（水变成水蒸气）；气 $\rightarrow$ 液：\textbf{液化}（水蒸气变成水，也叫凝结）。
\item 固 $\rightarrow$ 气：\textbf{升华}（樟脑丸变小、冻衣服变干）；气 $\rightarrow$ 固：\textbf{凝华}（窗玻璃上的冰花，是屋里的水蒸气直接冻成的）。
\end{itemize}

六个名字，背后是\textbf{同一条规律}：给粒子能量，它们就挣脱束缚、走向更自由的状态；粒子放出能量，就被束缚得更紧。熔化、汽化、升华都要\textbf{吸热}；凝固、液化、凝华都要\textbf{放热}。冰箱里冻冰块之所以会让冰箱耗电发热，就是因为水凝固时放出的能量必须被搬走。

有两个细节值得单独说一说。第一，汽化有两种方式：\textbf{蒸发}在任何温度下都发生——液面上总有一些跑得特别快的分子，能挣脱邻居的拉扯逃出去，所以湿衣服在阴雨天也能慢慢晾干；\textbf{沸腾}则要到特定温度，液体内部都咕嘟咕嘟冒出气泡才算。注意蒸发带走的是“跑得最快”的那批分子，留下来的分子平均就跑得慢了——所以蒸发会让剩下的液体\textbf{变凉}。你浑身是汗时扇扇子觉得凉快，发烧时医生用稀释的酒精擦皮肤降温，游泳上岸后风一吹直打哆嗦，都是同一条规律在起作用：逃走的分子把能量一起带走了。

第二，熔化、沸腾进行的时候，你继续加热，温度却\textbf{不再上升}——加进去的能量没有用来让粒子跑得更快，而是用来“拆散”粒子之间的束缚。这部分能量去哪儿了、为什么非花不可，第 27、28 章会用能量和熵的语言把账算清楚；而粒子之间到底靠什么互相拉住，是第 22 章的主角。

\begin{shiyan}
\textbf{观察“白气”从哪里来}（观察类实验，请家长操作热水，你在旁边看）

材料：一个烧水壶或带嘴的锅、一把手电筒（可选）、室温下的一杯水。

步骤：（1）请家长把水烧开，壶嘴开始冒白气后，关掉火源但保持观察。（2）仔细看壶嘴：白气是从壶嘴口开始的，还是离开壶嘴一小段距离才出现的？（3）用手电筒从侧面照白气，再照照壶嘴口那一段“看起来什么都没有”的透明区域，比较两者。（4）把一杯冷水举到白气上方（别碰到蒸汽，会烫伤），观察杯底。

观察点：壶嘴口那段透明区域里到底有没有东西？白气离壶嘴多远才开始出现？杯底发生了什么？

安全提示：沸水和水蒸气都能造成严重烫伤，绝对不要用手去摸壶嘴、蒸汽或“白气”，全程保持一臂距离。
\end{shiyan}

如果你观察仔细，会发现：壶嘴口那段\textbf{透明}的区域里才是真正的水蒸气——水蒸气是无色透明的气体，你的眼睛根本看不见它。白气是水蒸气冲出壶嘴后遇冷、重新凝结成的无数\textbf{小液滴}，本质上是一团极细的雾。所以“白气”不是气，是液体。这个发现等下回答开场问题时还要用。

\begin{qiaoliang}
液体变成气体，体积会涨多少？来算一笔真实的账。取 \ud{18}{mL} 水（大约一汤匙），它的质量是 \ud{18}{g}，正好是 $6.02\times 10^{23}$ 个水分子（这个巨大的计数单位是第 5 章的主题，这里先借用结论）。这么多分子如果全部变成 $100\,^{\circ}\mathrm{C}$、常压下的水蒸气，大约占 \ud{30}{L}——一个双肩背包那么大。也就是说，体积膨胀了约 $30000\div 18\approx 1700$ 倍。粒子本身并没有变大，是粒子之间的平均距离拉开了约 $\sqrt[3]{1700}\approx 12$ 倍。反过来，把气体压缩成液体，体积能缩小近千倍——这就是为什么一罐小小的液化气能供炉灶烧很久。密封容器里液体受热汽化、体积想涨一千多倍却涨不了，压强就会急剧升高，所以灌装气体的钢瓶都标着“严禁曝晒、严禁靠近火源”，这不是吓唬人。
\end{qiaoliang}

\section{给状态变化配上符号}

学到这里，可以请出符号了。化学里用括号里的小写字母标注状态：固态 s、液态 l、气态 g、溶解在水中 aq。于是冰化成水写作：

\[\chem{H_2O(s)} \rightarrow \chem{H_2O(l)}\]

每种物质都有自己标志性的转变温度，叫\textbf{熔点}和\textbf{沸点}。它们是物质的“指纹”之一，第 3 章讲的蒸馏、结晶等分离手段，靠的就是不同物质这些“指纹”的差异。下面是几种物质的熔点和沸点：

\begin{table}[htbp]
\centering
\begin{tabular}{lcc}
\toprule
物质 & 熔点（$^{\circ}\mathrm{C}$） & 沸点（$^{\circ}\mathrm{C}$） \\
\midrule
水 & 0 & 100 \\
酒精（乙醇） & $-114$ & 78 \\
氮气 & $-210$ & $-196$ \\
氧气 & $-218$ & $-183$ \\
铁 & 1538 & 2861 \\
\bottomrule
\end{tabular}
\end{table}

看一眼这张表，你就能读懂很多日常安排：北方冬天的温度计里装酒精而不装水，因为水在 $0\,^{\circ}\mathrm{C}$ 就凝固了，酒精到 $-114\,^{\circ}\mathrm{C}$ 还是液体；空气先被降温到 $-200\,^{\circ}\mathrm{C}$ 左右变成液体，再慢慢升温，沸点不同的氮气和氧气就能先后“蒸发”出来分开——工业上制取氧气用的正是这套办法。

把一杯冰持续加热、记录温度随时间的变化，会画出这样一条折线：

\begin{center}
\begin{tikzpicture}[scale=0.9]
\draw[->] (0,0) -- (9.2,0) node[right] {时间};
\draw[->] (0,0) -- (0,4.2) node[above] {温度};
\draw[thick] (0,0.4) -- (1.6,1.6) -- (3.4,1.6) -- (5.6,3.4) -- (8.2,3.4);
\draw[dashed] (0,1.6) -- (3.4,1.6);
\draw[dashed] (0,3.4) -- (8.2,3.4);
\node[left] at (0,1.6) {熔点};
\node[left] at (0,3.4) {沸点};
\node[below] at (0.8,0.9) {冰升温};
\node[below] at (2.5,1.6) {熔化中};
\node[below] at (4.5,2.4) {水升温};
\node[below] at (6.9,3.4) {沸腾中};
\end{tikzpicture}
\end{center}

两段平台对应熔化和沸腾：热量一直在加，温度却原地不动。这幅图是人为的示意，平台的宽窄、折线的斜率都不代表精确数据，但它把“能量先去拆束缚、再提高温度”这件事画得很清楚。

\section{状态变化不是化学变化}

现在回到第 1 章立下的规矩：物理变化和化学变化有什么区别？状态变化是最好的试金石。

水在冰、水、水蒸气之间怎么折腾，水分子始终是 \chem{H_2O}——两个氢原子和一个氧原子的组合没有拆开，没有重组，只是排列和运动的松紧变了。所以状态变化是\textbf{物理变化}：粒子本身没变，变的是粒子的“队形”。判断的硬标准不是“看起来变了没有”，而是“粒子有没有重新组合”。你有反证手段：把白气冷凝回水，它还是原来那杯水，烧干了也没剩下新东西。

反过来，如果用通电的办法把水分子拆开，变成氢气和氧气（第 33 章会讲），那才是化学变化——粒子被拆散重组成新分子，冷凝、降温这些办法再也无法让它们变回水。状态变化改变的只是队形，化学变化改变的是队员。
\section{“永久气体”是怎样被推翻的}

\begin{zhengju}
今天你知道，任何气体冷却到足够低的温度都能变成液体。但在 19 世纪初，这是一件没人敢打包票的事。当时的化学家能轻松液化氯气、氨气，可氧气、氮气、氢气不管加多大压强都“顽固不化”，于是教科书里出现了“永久气体”这个类别——被认为永远不会液化的气体。

转机来自一连串“意外”和一次关键的清醒。1823 年，法拉第在一根密封玻璃管里加热一种含氯的化合物，本想做别的实验，却发现管壁上凝出了油状的液体——他无意中同时给氯气加了压又降了温，氯气液化了。他顺着这条线索，用“加压 + 降温”液化了十几种气体，但氧气、氮气、氢气始终失败。

真正点破迷局的是爱尔兰物理学家安德鲁斯。1869 年，他系统研究二氧化碳在不同温度、压强下的表现，发现了一个惊人的现象：只要温度高于 $31\,^{\circ}\mathrm{C}$，无论加多大的压强，二氧化碳都不会液化，气体和液体的界线干脆消失了。他由此提出“临界温度”的概念：每种气体都有自己的临界温度，高于它，液化根本无从谈起。所谓的“永久气体”，其实只是临界温度特别低、当时的制冷技术够不着而已。1877 年，氧气和氮气几乎同时被液化；1908 年，荷兰物理学家昂内斯把最难啃的氦气也液化了，温度低到 $-269\,^{\circ}\mathrm{C}$。“永久气体”这个类别从此从教科书里划掉了。

这个故事里最值得带走的不是年份和人名，而是那个思维动作：失败没有催生“这些气体天生特殊”的结论，而是逼出一个更深刻的问题——是不是还缺了某个条件？一个被实验反复打脸的概念，就该让位给更好的概念。安德鲁斯发现的“液气界线消失”的现象，等下还会再登场。
\end{zhengju}

\section{三态之外的世界}

“固体、液体、气体”是一张很好用的地图，但它画的是常温常压附近的世界。把温度、压强推到极端，或者物质本身的结构复杂一点，地图的边界就开始模糊。下面这些“状态”，每一个都真实地出现在你的生活或新闻里。

\textbf{等离子体}。把气体加热到几千度以上，粒子撞得越来越狠，狠到把原子里的电子都撞了出来。于是气体不再由完整的中性粒子组成，而变成一锅自由电子和带电离子的“汤”——这就是等离子体，常被称为“物质的第四态”。因为带电，它会乖乖听从电场和磁场的指挥，这是它和普通气体最大的不同。闪电划过空气的通道、霓虹灯管里的辉光、电焊的弧光，都是等离子体；太阳和几乎所有恒星，整个就是巨大的等离子体球。按质量算，我们能看见的宇宙物质里，绝大多数处于这种状态——“三态”反而才是宇宙里的少数派。

\textbf{液晶}。用固体加热的方式熔化某些物质，会得到一种怪东西：它能像液体一样流动，内部分子却不像普通液体那样东倒西歪，而是保持着某种整齐的指向，有点晶体的派头。这种“流动着的有序”就是液晶。你的手机屏幕里就住着一层液晶：加不加电压，分子的指向会改变，光线的通过方式随之改变，像素就此点亮或熄灭。它不固不液又亦固亦液，“三态”的两个盒子都装不下它。

\textbf{玻璃态}。玻璃摸起来是地地道道的固体，可如果把镜头推进到粒子层面，它的粒子排列并不像晶体那样整齐，而是乱糟糟的，更像液体某一刻的“抓拍”——只不过这些粒子被冻得再也挪不动了。这种“结构像液体、行为像固体”的状态叫玻璃态。窗户玻璃、玻璃杯、玻璃纤维，还有冰箱里的果冻、硬化的树脂，都属于这一类。第 24 章讲晶体和新材料时，还会回来仔细看它。

\begin{wuqu}
“玻璃其实是流动得极慢的液体——你看欧洲老教堂的窗玻璃，都是下面厚上面薄，就是几百年流下来的证据。”这个说法流传很广，但它不成立。玻璃的流动需要粒子整体挪动位置，而在室温下，玻璃粒子“挪一步”所需的时间比宇宙的年龄还要长得多，几百年里流出的厚度连仪器的测量极限都够不着。那老窗玻璃为什么下厚上薄？因为老式平板玻璃本来就是吹制或甩制成的圆片，厚度天然不均匀，而装玻璃的工匠习惯把重的一边朝下装——事实上也发现了装反的、上厚下薄的老窗户。玻璃态是固体的一个真实类别，不是“慢到看不见的液体”。
\end{wuqu}

\textbf{超临界流体}。还记得安德鲁斯的发现吗：超过某个温度和压强（对二氧化碳是 $31\,^{\circ}\mathrm{C}$、约 73 个大气压），液体和气体的界线就消失了。在这个区域里，物质既不是液体也不是气体，而是一种两者性质兼有的“超临界流体”：它能像气体一样钻进缝隙，又能像液体一样溶解不少物质。这个状态听起来遥远，其实可能就在你的茶杯里——许多速溶咖啡和茶叶是用超临界二氧化碳脱掉咖啡因的，咖啡因被溶解带走，香味物质大多留下，二氧化碳一减压就全部跑光，不留任何残留。

所以本章标题说的“远远不止三态”，不是修辞。三态模型在你身边这个温度压强范围内非常好用，水、空气、金属的日常表现都能用它解释；但一旦温度过高、压强过大，或者物质的粒子排列有讲究，这张地图就得换更细的版本。科学模型没有“对错”之分那么简单，只有“在多大的范围内好用”。

\section{回到那袋消失的“冰”}

现在可以回答开场的三个问题了。

干冰是固态的二氧化碳，温度约 $-78.5\,^{\circ}\mathrm{C}$。它消失得干干净净，是因为在普通大气压下，二氧化碳\textbf{根本没有液态这一阶段}——固体受热后不经过液体，直接升华成气体跑掉了。为什么二氧化碳这么“急”，而水却老老实实地走“冰—水—水蒸气”的路线？这与每种物质在三态图上的版图有关，挑战框里会给出更完整的回答。

那团“白烟”也不是二氧化碳——二氧化碳和水蒸气一样，是无色透明的气体，你的眼睛看不见它。白烟是干冰把周围的空气急剧冷却，空气里的水蒸气凝结成的无数小液滴，和壶嘴那团白气是同一类东西。舞台上的“云雾”效果、冰镇饮料罐外壁的水珠，全都是这个道理：不是罐里的水渗出来了，是空气中的水蒸气在冷表面上缴械投降。

现在你也能想明白说明书上那句警告了：$79\,^{\circ}\mathrm{C}$ 左右的温差，意味着干冰接触皮肤的瞬间会让皮肤组织里的水急速凝固——徒手抓干冰，造成的伤害和严重烫伤是一个级别。看不见的危险，往往比看得见的更值得敬畏。

\section{从冻干咖啡到恒星}

状态变化这套规律，早就不只是课本里的名词，它在很多行业的流水线上默默干活。

食品工业里有一种“冷冻干燥”技术：先把食物冻透，再抽走空气降低压强，让冰直接升华跑掉。因为不经过液态，食物的细胞结构不被泡坏，营养和风味保留得最好——冻干咖啡、冻干水果、航天员的饭菜都靠它，一些怕热的疫苗和药品也用同样的思路保存。农业上，种子入库前要控制含水量，本质也是在管理水分在固、液、气之间的去留。

医学和工业离不开液态气体：医院的液氧、焊接用的液氮、保存细胞和组织的液氮罐，都是把气体液化后“浓缩”几百上千倍，便于储存和运输。等离子体则藏在日光灯管、芯片工厂和核聚变实验装置里——科学家正在尝试用磁场约束上亿度的等离子体，模仿太阳发电，这是人类能源的终极赌局之一（第 8 章会讲它背后的核反应）。

而所有这些状态差异——为什么水在 $0\,^{\circ}\mathrm{C}$ 结冰、酒精在 $-114\,^{\circ}\mathrm{C}$ 才结冰，为什么二氧化碳在常压下连液态都懒得出现，为什么水结冰后不但不收缩、反而膨胀让冰浮在水面上——最终都指向同一个幕后角色：\textbf{分子与分子之间那看不见的吸引力}。它决定了粒子被拉得多紧、要多少能量才能挣脱。第 22 章会把这只“看不见的手”拆给你看，第 23 章则专门审判水这个处处反常的“怪胎”。

\begin{tiaozhan}
（本框涉及相图，跳过不影响主线。）把温度和压强分别画在横轴和纵轴上，可以画出一张“状态地图”，叫相图：图上的曲线把版图分成固、液、气三块区域，曲线本身是“两态恰好能共存”的条件。这张图上有一个独一无二的点，叫\textbf{三相点}——在这个温度和压强下，固、液、气三态可以同时存在。水的三相点在约 $0.01\,^{\circ}\mathrm{C}$、$0.006$ 个大气压处，远低于正常大气压，所以你在常压下能看到完整的“冰—水—水蒸气”三部曲。而二氧化碳的三相点在 $-56.6\,^{\circ}\mathrm{C}$、约 $5.1$ 个大气压处——液态二氧化碳只存在于 5.1 个大气压\textbf{以上}的区域。普通大气压够不着它的液态版图，于是干冰受热只能直接从固态跨进气态，这就是它“不流汗”的真正原因。灭火器和钢瓶里的二氧化碳之所以是液体，正是因为瓶内压强远远高于这个门槛。
\end{tiaozhan}

\section*{练一练}

\begin{sikaoti}
\item 画三幅粒子示意图，分别表示同一种物质的固态、液态和气态。用点的疏密表示间距，用小短线表示运动方式，并在图旁注明：这是帮助思考的表示法，粒子既不是点，也不长这样。
\item 冬天结冰的衣服晾在室外也能慢慢变干。用粒子图像说明：冰没有熔化，水分子是怎样“逃走”的？这个过程叫什么？为什么阳光好的天气干得更快？
\item 高压锅煮饭更快。结合“压强—沸点”的关系（压强越大，沸点越高），说明高压锅里发生了什么，以及它为什么能把食物煮得更烂。再想一步：在海拔很高的高原上，不用高压锅为什么煮不熟饭？
\item 估算：液氮的沸点是 $-196\,^{\circ}\mathrm{C}$，\ud{1}{L} 液氮在室温下全部气化，体积大约膨胀多少倍？（提示：可以参照正文中水变水蒸气的算法，结论约为几百倍量级即可。）根据你的估算，解释为什么保存液氮的容器绝不能完全密封。
\item 设计一个观察或实验，向一位坚信“壶嘴的白气就是水蒸气”的朋友证明：白气是小液滴，而真正的水蒸气是看不见的。写出你的做法和预期现象。
\item 一杯水放在桌上，几天后消失了。画一幅物质流图：从杯子出发，画出水分子的去向，标出沿途经过的状态和推动每一步的条件。如果此时房间里有人关紧门窗并放了一盆冰块，你的图上会增加哪条路径？
\end{sikaoti}
